% SGA 3, Expose IV, English reconstruction.
% Scope: Section 4.5 opening through Proposition 4.5.1 and the definition of
% the induced topology on C/S, stopping before Proposition 4.5.2.
% Authority: Polo--Gille Exp4-13oct24.pdf, component-local pages 43--45
% (pp. 43 and 45 partial; combined-reader pages 231--233).
% The external English Markdown was used only as a comparison.

\setcounter{footnote}{39}

\SGARefTarget{sga3:iv:subsection:4:5}\subsection*{4.5. 比典范拓扑更粗的拓扑的情形}
\addcontentsline{toc}{subsection}{\SGARefLink{sga3:iv:subsection:4:5}{4.5}. 比典范拓扑更粗的拓扑}

由 \SGARefLink{sga3:iv:definition:4:3:6}{4.3.6} 和
\SGARefLink{sga3:iv:definition:4:3:8}{4.3.8},
关于 \(\mathscr C\) 上的一个拓扑 \(\mathcal T\) 的以下条件等价:
\begin{enumerate}
  \item[(i)]\SGARefTarget{sga3:iv:subsection:4:5:item:i:01}
  \(\mathcal T\) 比 \(\mathscr C\) 的典范拓扑更粗;
  \item[(ii)]\SGARefTarget{sga3:iv:subsection:4:5:item:ii:01}
  每个可表预层都是关于 \(\mathcal T\) 的层;
  \item[(iii)]\SGARefTarget{sga3:iv:subsection:4:5:item:iii:01}
  每个关于 \(\mathcal T\) 的覆盖筛都具有普遍有效满态射性.
\end{enumerate}
若 \(\mathcal T\) 由一个预拓扑 \(S\mapsto\mathscr R(S)\) 定义,
则这些条件还等价于:
\begin{enumerate}
  \item[(iv)]\SGARefTarget{sga3:iv:subsection:4:5:item:iv:01}
  属于 \(\mathscr R(S)\) 的每个族都是普遍有效满态射族.
\end{enumerate}

当这些条件成立时, 典范函子
\(\mathscr C\to\widehat{\mathscr C}\) 通过一个函子
\[
  j_{\mathscr C}=j:\mathscr C\longrightarrow\widetilde{\mathscr C};
\]
分解; 也记 \(j(S)=\widetilde S\).\footnote{编者注: 还将记
\(\mathbf h_S=\widehat S\); 参见 \SGARefLink{sga3:iv:corollary:4:5:4}{4.5.4}
中的第一个交换图.}

\medskip
\noindent\SGARefTarget{sga3:iv:proposition:4:5:1}\textbf{命题 4.5.1.}
\textit{函子 \(j:\mathscr C\to\widetilde{\mathscr C}\) 是全忠实的,
且与任意逆极限可交换. 特别地, 它左正合,
因而保持由有限逆极限定义的代数结构.}

这立即由交换的分解图
\[
\begin{array}{ccc}
 &\mathscr C&\\[-0.5ex]
 {\scriptstyle j}\swarrow&&\searrow\\[-0.5ex]
 \widetilde{\mathscr C}&\xrightarrow{\ i\ }&\widehat{\mathscr C}
\end{array}
\]
及 \SGARefLink{sga3:iv:theorem:4:4:1:item:i:01}{4.4.1(i)} 得到.

在给出 \(j\) 的更多性质之前, 必须先定义范畴 \(\mathscr C/S\)
上的诱导拓扑. 此处不再假设给定的拓扑必定比典范拓扑更粗.
\(\mathscr C\) 中对象 \(T\) 的一个筛, 连同一个态射 \(T\to S\),
自然地定义了 \(\mathscr C/S\) 中 \(T\) 的一个筛,
因为 \(T\) 的筛的\SGARefLink{sga3:iv:definition:4:1:2}{定义}
只依赖于范畴
\[
  \mathscr C/T=(\mathscr C/S)/T.
\]
例如, 若这个筛由一个族 \(\{T_i\to T\}\) 定义,
则它在 \(\mathscr C/S\) 中的像由同一个族定义,
只是现在将此族视为 \(\mathscr C/S\) 中的态射族.

规定覆盖筛的这些像为覆盖筛, 便在 \(\mathscr C/S\) 上定义了一个拓扑,
称为给定拓扑的\emph{诱导拓扑 (induced topology)}.
用 \SGARefLink{sga3:iv:bib:as}{[AS]}, \SGARefLink{sga3:iv:proposition:2:3}{2.3}
的术语来说, 这是 \(\mathscr C/S\) 上使典范函子
\[
  i_S:\mathscr C/S\longrightarrow\mathscr C
\]
成为余态射的最粗拓扑.\footnote{编者注: 即对 \(\mathscr C/S\)
的每个对象 \(T\to S\), \(i_S(T\to S)=T\) 的每个覆盖筛
\(R'\), 视为对象 \((T\to S)\in\operatorname{Ob}(\mathscr C/S)\)
的筛时仍为覆盖筛.}
按定义, 等同
\[
  (\mathscr C/S)/T=\mathscr C/T
\]
保持这些拓扑.

% Translation coverage: entire assigned source component, without omissions.
% Statistics: source lines=83; target lines=87; inline math=31; display math=5; unresolved translation TODOs=0.
